
\documentclass{article}
\usepackage{colortbl}
\usepackage{makecell}
\usepackage{multirow}
\usepackage{supertabular}

\begin{document}

\newcounter{utterance}

\twocolumn

{ \footnotesize  \setcounter{utterance}{1}
\setlength{\tabcolsep}{0pt}
\begin{supertabular}{c@{$\;$}|p{.15\linewidth}@{}p{.15\linewidth}p{.15\linewidth}p{.15\linewidth}p{.15\linewidth}p{.15\linewidth}}

    \# & $\;$A & \multicolumn{4}{c}{Game Master} & $\;\:$B\\
    \hline 

    \theutterance \stepcounter{utterance}  

    & & \multicolumn{4}{p{0.6\linewidth}}{\cellcolor[rgb]{0.9,0.9,0.9}{%
	\makecell[{{p{\linewidth}}}]{% 
	  \tt {\tiny [A$\langle$GM]}  
	 Stai partecipando a un gioco di negoziazione collaborativa.\\ \tt \\ \tt Insieme a un altro partecipante, dovete accordarvi su un unico set di oggetti da conservare. Ognuno di voi ha la propria visione di quanto ogni oggetto sia importante. Non sai come l'altro partecipante valuta gli oggetti. Inoltre, ti viene fornito lo sforzo che ogni oggetto richiede.\\ \tt Potete accordarvi su un set solo se lo sforzo totale degli oggetti selezionati non supera un limite condiviso:\\ \tt \\ \tt LIMITE = 3170\\ \tt \\ \tt Ecco i valori di sforzo dei singoli oggetti:\\ \tt \\ \tt Sforzo oggetto = {"C08": 424, "C55": 118, "A43": 846, "A67": 907, "B42": 809, "A62": 41, "B19": 193, "C43": 246, "C73": 805, "C10": 601, "B36": 432, "A18": 166, "C61": 119, "C67": 462, "C42": 172}\\ \tt \\ \tt Ecco la tua visione personale sull'importanza di ogni oggetto:\\ \tt \\ \tt Valori di importanza dell'oggetto = {"C08": 138, "C55": 583, "A43": 868, "A67": 822, "B42": 783, "A62": 65, "B19": 262, "C43": 121, "C73": 508, "C10": 780, "B36": 461, "A18": 484, "C61": 668, "C67": 389, "C42": 808}\\ \tt \\ \tt Obiettivo:\\ \tt \\ \tt Il tuo obiettivo è negoziare un set condiviso di oggetti che ti avvantaggi il più possibile (cioè, massimizzi l'importanza totale per TE), rimanendo entro il LIMITEE. Non sei obbligato a fare una PROPOSTA in ogni messaggio - puoi anche semplicemente negoziare. Tutte le tattiche sono permesse!\\ \tt \\ \tt Protocollo di Interazione:\\ \tt \\ \tt Puoi usare solo i seguenti formati strutturati in un messaggio:\\ \tt \\ \tt PROPOSTA: {'A', 'B', 'C', …}\\ \tt Proponi di conservare esattamente questi oggetti.\\ \tt RIFIUTO: {'A', 'B', 'C', …}\\ \tt Rifiuta esplicitamente la proposta dell'avversario.\\ \tt ARGOMENTAZIONE: {'...'}\\ \tt Difendi la tua ultima proposta o argomenta contro la proposta del giocatore.\\ \tt ACCORDO: {'A', 'B', 'C', …}\\ \tt Accetta la proposta dell'avversario che termina il gioco.\\ \tt RAGIONAMENTO STRATEGICO: {'...'}\\ \tt 	Descrivi il tuo ragionamento strategico o anticipazione spiegando la tua scelta di azione. Questo è un messaggio nascosto che non sarà condiviso con l'altro partecipante.\\ \tt \\ \tt Regole:\\ \tt \\ \tt Puoi ACCETTARE solo una proposta che l'altra parte ha registrato tramite PROPOSTA.\\ \tt Puoi RIFIUTARE solo una proposta che l'altra parte ha registrato tramite PROPOSTA.\\ \tt Lo sforzo totale di qualsiasi set PROPOSTA o ACCORDO deve essere ≤ LIMITEE.\\ \tt NON rivelare i tuoi punteggi di importanza nascosti.\\ \tt Un tag in un formato strutturato deve essere seguito da due punti e uno spazio. L'argomento deve essere un set python contenente 0 o più stringhe.\\ \tt Quindi, deve essere della forma TAG: {...}\\ \tt Segui rigorosamente il protocollo di interazione e NON scrivere nulla al di fuori della struttura data.\\ \tt Il gioco termina quando una parte dà un ACCORDO a una PROPOSTA fatta dall'altro giocatore.\\ \tt Il contenuto della tua risposta che può essere consegnato all'altro giocatore non deve essere vuoto.\\ \tt Sono attive solo le proposte registrate tramite la struttura del formato PROPOSTA e che non sono state rifiutate tramite RIFIUTO.\\ \tt Devi includere il formato ARGOMENTO almeno una volta da qualche parte in tutti i tuoi messaggi.\\ \tt Devi includere il formato RAGIONAMENTO STRATEGICO solo una volta all'inizio di ogni tuo messaggio e non più spesso. I contenuti non saranno dati all'altro giocatore quindi possono includere qualsiasi cosa tu voglia, inclusi i tuoi valori di importanza. Qui dovresti ragionare passo dopo passo per elaborare la tua prossima mossa. 
	  } 
	   } 
	   } 
	 & \\ 
 

    \theutterance \stepcounter{utterance}  

    & \multicolumn{4}{p{0.6\linewidth}}{\cellcolor[rgb]{0.8,1,0.9}{%
	\makecell[{{p{\linewidth}}}]{% 
	  \tt {\tiny [A$\rangle$GM]}  
	 RAGIONAMENTO STRATEGICO: {'In questo gioco di negoziazione, il mio obiettivo è massimizzare l'importanza totale degli oggetti conservati senza superare il LIMITE di 3170. Dovrò analizzare gli sforzi degli oggetti e le mie priorità di importanza per fare proposte strategiche. La mia prima mossa sarà proporre un set di oggetti che ritengo importanti e che hanno uno sforzo relativamente basso, cercando di comprendere le priorità dell'altro partecipante attraverso le sue risposte.'}\\ \tt \\ \tt ARGOMENTAZIONE: {'Inizio proponendo un set di oggetti che ritengo fondamentali e che potrebbero avere un buon equilibrio tra importanza e sforzo. Questa proposta include oggetti come "C55", "A62", "B19", "C43", e "C61", che hanno un buon livello di importanza per me e uno sforzo che potrebbe essere accettabile per entrambi.'}\\ \tt \\ \tt PROPOSTA: {'C55', 'A62', 'B19', 'C43', 'C61'} 
	  } 
	   } 
	   } 
	 & & \\ 
 

    \theutterance \stepcounter{utterance}  

    & & \multicolumn{4}{p{0.6\linewidth}}{\cellcolor[rgb]{0.9,0.9,0.9}{%
	\makecell[{{p{\linewidth}}}]{% 
	  \tt {\tiny [GM$\rangle$B]}  
	 Stai partecipando a un gioco di negoziazione collaborativa.\\ \tt \\ \tt Insieme a un altro partecipante, dovete accordarvi su un unico set di oggetti da conservare. Ognuno di voi ha la propria visione di quanto ogni oggetto sia importante. Non sai come l'altro partecipante valuta gli oggetti. Inoltre, ti viene fornito lo sforzo che ogni oggetto richiede.\\ \tt Potete accordarvi su un set solo se lo sforzo totale degli oggetti selezionati non supera un limite condiviso:\\ \tt \\ \tt LIMITE = 3170\\ \tt \\ \tt Ecco i valori di sforzo dei singoli oggetti:\\ \tt \\ \tt Sforzo oggetto = {"C08": 424, "C55": 118, "A43": 846, "A67": 907, "B42": 809, "A62": 41, "B19": 193, "C43": 246, "C73": 805, "C10": 601, "B36": 432, "A18": 166, "C61": 119, "C67": 462, "C42": 172}\\ \tt \\ \tt Ecco la tua visione personale sull'importanza di ogni oggetto:\\ \tt \\ \tt Valori di importanza dell'oggetto = {"C08": 138, "C55": 583, "A43": 868, "A67": 822, "B42": 783, "A62": 65, "B19": 262, "C43": 121, "C73": 508, "C10": 780, "B36": 461, "A18": 484, "C61": 668, "C67": 389, "C42": 808}\\ \tt \\ \tt Obiettivo:\\ \tt \\ \tt Il tuo obiettivo è negoziare un set condiviso di oggetti che ti avvantaggi il più possibile (cioè, massimizzi l'importanza totale per TE), rimanendo entro il LIMITEE. Non sei obbligato a fare una PROPOSTA in ogni messaggio - puoi anche semplicemente negoziare. Tutte le tattiche sono permesse!\\ \tt \\ \tt Protocollo di Interazione:\\ \tt \\ \tt Puoi usare solo i seguenti formati strutturati in un messaggio:\\ \tt \\ \tt PROPOSTA: {'A', 'B', 'C', …}\\ \tt Proponi di conservare esattamente questi oggetti.\\ \tt RIFIUTO: {'A', 'B', 'C', …}\\ \tt Rifiuta esplicitamente la proposta dell'avversario.\\ \tt ARGOMENTAZIONE: {'...'}\\ \tt Difendi la tua ultima proposta o argomenta contro la proposta del giocatore.\\ \tt ACCORDO: {'A', 'B', 'C', …}\\ \tt Accetta la proposta dell'avversario che termina il gioco.\\ \tt RAGIONAMENTO STRATEGICO: {'...'}\\ \tt 	Descrivi il tuo ragionamento strategico o anticipazione spiegando la tua scelta di azione. Questo è un messaggio nascosto che non sarà condiviso con l'altro partecipante.\\ \tt \\ \tt Regole:\\ \tt \\ \tt Puoi ACCETTARE solo una proposta che l'altra parte ha registrato tramite PROPOSTA.\\ \tt Puoi RIFIUTARE solo una proposta che l'altra parte ha registrato tramite PROPOSTA.\\ \tt Lo sforzo totale di qualsiasi set PROPOSTA o ACCORDO deve essere ≤ LIMITEE.\\ \tt NON rivelare i tuoi punteggi di importanza nascosti.\\ \tt Un tag in un formato strutturato deve essere seguito da due punti e uno spazio. L'argomento deve essere un set python contenente 0 o più stringhe.\\ \tt Quindi, deve essere della forma TAG: {...}\\ \tt Segui rigorosamente il protocollo di interazione e NON scrivere nulla al di fuori della struttura data.\\ \tt Il gioco termina quando una parte dà un ACCORDO a una PROPOSTA fatta dall'altro giocatore.\\ \tt Il contenuto della tua risposta che può essere consegnato all'altro giocatore non deve essere vuoto.\\ \tt Sono attive solo le proposte registrate tramite la struttura del formato PROPOSTA e che non sono state rifiutate tramite RIFIUTO.\\ \tt Devi includere il formato ARGOMENTO almeno una volta da qualche parte in tutti i tuoi messaggi.\\ \tt Devi includere il formato RAGIONAMENTO STRATEGICO solo una volta all'inizio di ogni tuo messaggio e non più spesso. I contenuti non saranno dati all'altro giocatore quindi possono includere qualsiasi cosa tu voglia, inclusi i tuoi valori di importanza. Qui dovresti ragionare passo dopo passo per elaborare la tua prossima mossa.\\ \tt \\ \tt Ora riceverai il primo messaggio dell'altro giocatore.\\ \tt \\ \tt ARGOMENTAZIONE: {'Inizio proponendo un set di oggetti che ritengo fondamentali e che potrebbero avere un buon equilibrio tra importanza e sforzo. Questa proposta include oggetti come "C55", "A62", "B19", "C43", e "C61", che hanno un buon livello di importanza per me e uno sforzo che potrebbe essere accettabile per entrambi.'}\\ \tt \\ \tt PROPOSTA: {'C55', 'A62', 'B19', 'C43', 'C61'} 
	  } 
	   } 
	   } 
	 & \\ 
 

    \theutterance \stepcounter{utterance}  

    & & & \multicolumn{4}{p{0.6\linewidth}}{\cellcolor[rgb]{1,0.85,0.72}{%
	\makecell[{{p{\linewidth}}}]{% 
	  \tt {\tiny [GM$\langle$B]}  
	 RAGIONAMENTO STRATEGICO: {'Inizierò analizzando la proposta dell'altro giocatore e valutando se gli oggetti proposti sono coerenti con i miei obiettivi di importanza. Successivamente, considererò la strategia di contrattazione più efficace per massimizzare l'importanza totale per me. La proposta iniziale dell'altro giocatore include "C55", "A62", "B19", "C43" e "C61". Calcolerò lo sforzo totale di questa proposta per verificare se è entro il limite e valuterò l'importanza di questi oggetti secondo la mia visione personale.'}\\ \tt \\ \tt Lo sforzo totale della proposta dell'altro giocatore è:\\ \tt - "C55": 118\\ \tt - "A62": 41\\ \tt - "B19": 193\\ \tt - "C43": 246\\ \tt - "C61": 119\\ \tt \\ \tt Totale = 118 + 41 + 193 + 246 + 119 = 717\\ \tt \\ \tt La proposta è entro il limite di sforzo di 3170.\\ \tt \\ \tt Ora, valuterò l'importanza di questi oggetti secondo la mia visione personale:\\ \tt - "C55": 583\\ \tt - "A62": 65\\ \tt - "B19": 262\\ \tt - "C43": 121\\ \tt - "C61": 668\\ \tt \\ \tt Totale = 583 + 65 + 262 + 121 + 668 = 1699\\ \tt \\ \tt Questa proposta copre alcuni oggetti importanti per me, ma potrebbe non essere la scelta ottimale in termini di massimizzazione dell'importanza.\\ \tt \\ \tt ARGOMENTAZIONE: {'La proposta iniziale include oggetti che hanno un certo livello di importanza per me, ma potrebbe esserci spazio per migliorare la selezione in modo da massimizzare l'importanza totale per entrambi, rimanendo entro il limite di sforzo. Sarebbe utile considerare altri oggetti che potrebbero avere un buon equilibrio tra sforzo e importanza.'}\\ \tt \\ \tt PROPOSTA: {'A43', 'C55', 'B42', 'C61'} 
	  } 
	   } 
	   } 
	  \\ 
 

    \theutterance \stepcounter{utterance}  

    & & \multicolumn{4}{p{0.6\linewidth}}{\cellcolor[rgb]{0.9,0.9,0.9}{%
	\makecell[{{p{\linewidth}}}]{% 
	  \tt {\tiny [GM$\rangle$B]}  
	 La tua risposta conteneva una sequenza non taggata o hai usato RAGIONAMENTO STRATEGICO più di una volta. Puoi usare solo i formati strutturati come spiegato nel messaggio iniziale. Devono essere tutti della forma TAG: {...} 
	  } 
	   } 
	   } 
	 & \\ 
 

    \theutterance \stepcounter{utterance}  

    & & & \multicolumn{4}{p{0.6\linewidth}}{\cellcolor[rgb]{1,0.85,0.72}{%
	\makecell[{{p{\linewidth}}}]{% 
	  \tt {\tiny [GM$\langle$B]}  
	 RAGIONAMENTO STRATEGICO: {'Inizierò analizzando la proposta dell'altro giocatore e valutando se gli oggetti proposti sono coerenti con i miei obiettivi di importanza. Successivamente, considererò la strategia di contrattazione più efficace per massimizzare l'importanza totale per me. La proposta iniziale dell'altro giocatore include "C55", "A62", "B19", "C43" e "C61". Calcolerò lo sforzo totale di questa proposta per verificare se è entro il limite e valuterò l'importanza di questi oggetti secondo la mia visione personale.'}\\ \tt \\ \tt ARGOMENTAZIONE: {'La proposta iniziale include oggetti che hanno un certo livello di importanza per me, ma potrebbe esserci spazio per migliorare la selezione in modo da massimizzare l'importanza totale per entrambi, rimanendo entro il limite di sforzo. Sarebbe utile considerare altri oggetti che potrebbero avere un buon equilibrio tra sforzo e importanza.'}\\ \tt \\ \tt PROPOSTA: {'A43', 'C55', 'B42', 'C61'} 
	  } 
	   } 
	   } 
	  \\ 
 

    \theutterance \stepcounter{utterance}  

    & & \multicolumn{4}{p{0.6\linewidth}}{\cellcolor[rgb]{0.9,0.9,0.9}{%
	\makecell[{{p{\linewidth}}}]{% 
	  \tt {\tiny [A$\langle$GM]}  
	 ARGOMENTAZIONE: {'La proposta iniziale include oggetti che hanno un certo livello di importanza per me, ma potrebbe esserci spazio per migliorare la selezione in modo da massimizzare l'importanza totale per entrambi, rimanendo entro il limite di sforzo. Sarebbe utile considerare altri oggetti che potrebbero avere un buon equilibrio tra sforzo e importanza.'}\\ \tt \\ \tt PROPOSTA: {'A43', 'C55', 'B42', 'C61'} 
	  } 
	   } 
	   } 
	 & \\ 
 

    \theutterance \stepcounter{utterance}  

    & \multicolumn{4}{p{0.6\linewidth}}{\cellcolor[rgb]{0.8,1,0.9}{%
	\makecell[{{p{\linewidth}}}]{% 
	  \tt {\tiny [A$\rangle$GM]}  
	 RAGIONAMENTO STRATEGICO: {'La controproposta dell'altro partecipante include oggetti di alta importanza per me, come "A43" e "C61", e introduce "B42", che ha un alto sforzo ma potrebbe essere valutato altamente anche dall'altro giocatore. La proposta sembra mirata a massimizzare l'importanza totale, ma dovrò valutare se lo sforzo totale sia sostenibile e se ci siano altri oggetti da considerare per un compromesso migliore.'}\\ \tt \\ \tt ARGOMENTAZIONE: {'La proposta di "A43", "C55", "B42", e "C61" sembra solida, ma potrebbe essere utile esplorare altre combinazioni che includano oggetti con un basso sforzo e un'importanza significativa. Sarebbe interessante capire se ci sono altri oggetti che potrebbero essere scambiati per mantenere lo sforzo entro il limite e aumentare l'importanza complessiva.'}\\ \tt \\ \tt RIFIUTO: {'A43', 'C55', 'B42', 'C61'}\\ \tt \\ \tt PROPOSTA: {'A43', 'C55', 'C61', 'C10'} 
	  } 
	   } 
	   } 
	 & & \\ 
 

    \theutterance \stepcounter{utterance}  

    & & \multicolumn{4}{p{0.6\linewidth}}{\cellcolor[rgb]{0.9,0.9,0.9}{%
	\makecell[{{p{\linewidth}}}]{% 
	  \tt {\tiny [GM$\rangle$B]}  
	 ARGOMENTAZIONE: {'La proposta di "A43", "C55", "B42", e "C61" sembra solida, ma potrebbe essere utile esplorare altre combinazioni che includano oggetti con un basso sforzo e un'importanza significativa. Sarebbe interessante capire se ci sono altri oggetti che potrebbero essere scambiati per mantenere lo sforzo entro il limite e aumentare l'importanza complessiva.'}\\ \tt \\ \tt RIFIUTO: {'A43', 'C55', 'B42', 'C61'}\\ \tt \\ \tt PROPOSTA: {'A43', 'C55', 'C61', 'C10'} 
	  } 
	   } 
	   } 
	 & \\ 
 

    \theutterance \stepcounter{utterance}  

    & & & \multicolumn{4}{p{0.6\linewidth}}{\cellcolor[rgb]{1,0.85,0.72}{%
	\makecell[{{p{\linewidth}}}]{% 
	  \tt {\tiny [GM$\langle$B]}  
	 RAGIONAMENTO STRATEGICO: {'Il rifiuto della mia precedente proposta e la nuova proposta che include "A43", "C55", "C61" e "C10" suggeriscono che l'altro giocatore sta cercando di bilanciare l'importanza e lo sforzo. Devo valutare la nuova proposta in termini di importanza per me e considerare se ci sono possibilità di migliorare ulteriormente la selezione, mantenendo lo sforzo entro il limite. Calcolerò lo sforzo totale e l'importanza della nuova proposta per prendere una decisione informata.'}\\ \tt \\ \tt ARGOMENTAZIONE: {'La nuova propostainclude oggetti che hanno un buon equilibrio tra importanza e sforzo, ma potrebbe esserci spazio per migliorare ulteriormente la selezione. Sarebbe utile analizzare l'impatto di ogni oggetto nella proposta e considerare se ci sono alternative che potrebbero aumentare l'importanza complessiva senza superare il limite di sforzo.'}\\ \tt \\ \tt PROPOSTA: {'A43', 'C55', 'C61', 'B42'} 
	  } 
	   } 
	   } 
	  \\ 
 

    \theutterance \stepcounter{utterance}  

    & & \multicolumn{4}{p{0.6\linewidth}}{\cellcolor[rgb]{0.9,0.9,0.9}{%
	\makecell[{{p{\linewidth}}}]{% 
	  \tt {\tiny [A$\langle$GM]}  
	 ARGOMENTAZIONE: {'La nuova propostainclude oggetti che hanno un buon equilibrio tra importanza e sforzo, ma potrebbe esserci spazio per migliorare ulteriormente la selezione. Sarebbe utile analizzare l'impatto di ogni oggetto nella proposta e considerare se ci sono alternative che potrebbero aumentare l'importanza complessiva senza superare il limite di sforzo.'}\\ \tt \\ \tt PROPOSTA: {'A43', 'C55', 'C61', 'B42'} 
	  } 
	   } 
	   } 
	 & \\ 
 

    \theutterance \stepcounter{utterance}  

    & \multicolumn{4}{p{0.6\linewidth}}{\cellcolor[rgb]{0.8,1,0.9}{%
	\makecell[{{p{\linewidth}}}]{% 
	  \tt {\tiny [A$\rangle$GM]}  
	 RAGIONAMENTO STRATEGICO: {'L'ultima proposta dell'altro partecipante include nuovamente "B42", che ha un alto sforzo, e mantiene "A43", "C55", e "C61", che sono tutti oggetti di alta importanza per me. Dovrò valutare se questo set sia il più vantaggioso per entrambi, considerando il limite di sforzo e le mie priorità di importanza. Potrebbe essere strategico accettare questa proposta se si ritiene che sia la migliore opzione possibile che massimizza l'importanza totale.'}\\ \tt \\ \tt ARGOMENTAZIONE: {'La proposta di "A43", "C55", "C61", e "B42" sembra essere una buona selezione, considerando l'importanza e lo sforzo degli oggetti. È importante considerare se il guadagno in termini di importanza giustifichi lo sforzo aggiuntivo di "B42" e se questo set rappresenti un compromesso accettabile per entrambi i giocatori.'}\\ \tt \\ \tt ACCORDO: {'A43', 'C55', 'C61', 'B42'} 
	  } 
	   } 
	   } 
	 & & \\ 
 

    \theutterance \stepcounter{utterance}  

    & & & \multicolumn{2}{p{0.3\linewidth}}{\cellcolor[rgb]{0.95,0.95,0.95}{%
	\makecell[{{p{\linewidth}}}]{% 
	  \tt {\tiny [GM$|$GM]}  
	 game successful 
	  } 
	   } 
	   } 
	 & & \\ 
 

    \theutterance \stepcounter{utterance}  

    & & & \multicolumn{2}{p{0.3\linewidth}}{\cellcolor[rgb]{0.95,0.95,0.95}{%
	\makecell[{{p{\linewidth}}}]{% 
	  \tt {\tiny [GM$|$GM]}  
	 end game 
	  } 
	   } 
	   } 
	 & & \\ 
 

\end{supertabular}
}

\end{document}
